% 验证与校核、已知局限
% 本章文件由 \input 引入，不含导言区。
\section{验证与校核、已知局限}\label{sec:verification}

本章汇总 OPF 模块的全部验证证据与已知局限。前面各章（第~\ref{sec:acopfnative}~--~\ref{sec:threephase}~节）
给出各求解器的模型与算法，本章回答两个问题：\textbf{凭什么相信这些求解器的结果}，
以及\textbf{在哪些规模、工况与元件组合下结果尚未被充分验证}。写作口径与
《元件模型技术手册》一致：测试断言一律给出\texttt{文件:行号}，文档记录的实测数值
标注“文档实测”，被门控、被跳过、被降级为警告的条目如实列出，不以路线图代替现状。

\subsection{验证体系与证据分级}\label{subsec:ver:layers}

OPF 验证按四个层次组织（图~\ref{fig:ver:layers}）：

\begin{enumerate}
  \item \textbf{导数级}：Parity 建模层的解析等式 Jacobian、LCC 三变量一阶导数与局部
        Hessian、相域模型导数，全部与中心有限差分比对
        （\srcpath{tests/test\_opf\_solver\_backends.cpp}:274--308、310--422、424--598；
        \srcpath{tests/test\_three\_phase\_hybrid\_opf.cpp}:174--282）；
  \item \textbf{优化级}：收敛后检查 KKT 残差（\fld{max\_constraint\_violation}、
        \fld{max\_stationarity}）、目标值符号与量级、变量边界、实际生效的求解路径
        （\fld{solver\_path}）与线性代数后端（\fld{linear\_solver\_backend}）；
  \item \textbf{系统级}：把 OPF 调度写回系统副本运行独立 Newton 潮流回放（PF replay）、
        \srcpath{verify\_opf\_result} 物理审计、图归约/投影前后结果 round-trip；
  \item \textbf{外部基准}：手算闭式解、MATPOWER 标准算例、DSP（BPA）LCC 参考解、
        文档记录的参考目标值。
\end{enumerate}

\begin{figure}[H]
\centering
\begin{tikzpicture}[
  font=\small,
  box/.style={draw, rectangle, minimum width=10.6cm, minimum height=0.9cm,
              align=center, inner sep=2pt},
  lbl/.style={font=\footnotesize, align=left, text width=4.6cm},
  arr/.style={-{Stealth[length=2mm]}, thick}]
  \node[box] (l1) at (0,0) {导数级：解析 Jacobian/Hessian vs 有限差分};
  \node[box] (l2) at (0,1.25) {优化级：KKT 残差、目标值、边界、求解路径};
  \node[box] (l3) at (0,2.5) {系统级：PF 回放、结果审计、投影 round-trip};
  \node[box] (l4) at (0,3.75) {外部基准：手算解、MATPOWER、DSP/BPA、参考目标值};
  \draw[arr] (l1.north) -- (l2.south);
  \draw[arr] (l2.north) -- (l3.south);
  \draw[arr] (l3.north) -- (l4.south);
  \node[lbl, anchor=west] at (5.6,0) {微算例即可，不依赖外部工具};
  \node[lbl, anchor=west] at (5.6,1.25) {内置于默认 ctest 套件};
  \node[lbl, anchor=west] at (5.6,2.5) {跨模块一致性};
  \node[lbl, anchor=west] at (5.6,3.75) {部分为文档实测，非逐轮断言};
\end{tikzpicture}
\caption{OPF 验证的四层结构。下层失败会立即定位到导数或模型装配错误；
上层失败则指向求解器全局化或外部语义差异。}
\label{fig:ver:layers}
\end{figure}

证据按可信度分三档，本章表格与行文中严格区分：
\begin{itemize}
  \item \textbf{默认 ctest 断言}：注册于 \srcpath{tests/CMakeLists.txt}:31--34、109、115、
        185、189、282、288 的测试目标，每次构建均可复跑；
  \item \textbf{隐藏性能用例}：带 Catch2 隐藏标签 \fld{[.performance]} 的用例
        （如 \srcpath{tests/test\_opf\_solver\_backends.cpp}:914--943、945--1024，
        \srcpath{tests/test\_reactive\_power\_opt.cpp}:84--110），默认测试运行不执行，
        其结果不构成逐轮回归判据；
  \item \textbf{文档实测}：\srcpath{docs/large\_hybrid\_ipm\_diagnostics.md}、
        \srcpath{docs/reactive\_power\_optimization\_validation.md}、
        \srcpath{docs/lcc\_dat\_opf\_report/} 中记录的墙钟时间、迭代数与残差，
        是特定机器、特定构建下的一次性记录，引用时标“文档实测”。
\end{itemize}

单位制同全书：容差中 pu 为标幺值（$S_{\mathrm{base}}=100$~MVA），功率容差为
MW/MVAr 有名值，角度容差为度或 rad（随测试原文注明），成本币种未定。

\subsection{验证矩阵总表}\label{subsec:ver:matrix}

表~\ref{tab:ver:matrix} 汇总全部验证通道。各通道细节见后续小节；表中“容差/通过准则”
列均为测试断言或文档声明的原文数值。

\begin{footnotesize}
\begin{longtable}{@{}L{2.0cm}L{3.0cm}L{2.9cm}L{3.1cm}L{3.9cm}@{}}
\caption{OPF 验证矩阵总表}\label{tab:ver:matrix}\\
\toprule
\textbf{验证通道} & \textbf{算例（规模）} & \textbf{对照基准} & \textbf{容差/通过准则} & \textbf{出处（文件:行号）}\\
\midrule
\endfirsthead
\multicolumn{5}{@{}l}{\footnotesize（续表）}\\
\toprule
\textbf{验证通道} & \textbf{算例（规模）} & \textbf{对照基准} & \textbf{容差/通过准则} & \textbf{出处（文件:行号）}\\
\midrule
\endhead
\midrule
\multicolumn{5}{r@{}}{\footnotesize（接下页）}\\
\endfoot
\bottomrule
\endlastfoot
手算解析解 &
\fld{opf\_hand\_01}--\fld{04}（1--2 母线 JSON 微算例） &
README 闭式手算（merit order、支路拥塞、负荷削减、线性化成本） &
LP 精确；QP 参考目标 815.333 &
\srcpath{external\_data/opf\_hand\_cases/README.md}:21--296（未接入 ctest，GUI 手检）\\
MATPOWER 小算例 &
case5/9/14/30/39/57（$\le 57$ 节点） &
NativeAC 与 NativeQP 各自收敛 + Newton PF 回放 &
可行性 $10^{-6}$，PF 残差 $10^{-8}$ &
\srcpath{tests/test\_acopf\_dcopf\_crossval.cpp}:52--54、117--144、250--303；
\srcpath{tests/test\_matpower\_cases.cpp}:618--638\\
MATPOWER 大算例 &
case2383wp、case2869pegase（2869 节点）、case6515rte（6515 节点） &
Parity IPM 收敛/有限性；文档参考目标值 133999/386107 &
case2869：收敛且迭代 $<100$、稀疏后端；case6515rte：仅断言有限 &
\srcpath{tests/test\_opf\_solver\_backends.cpp}:1058--1071、1154--1196；
\srcpath{docs/large\_hybrid\_ipm\_diagnostics.md}:263--269（文档实测）\\
AC/DC/PF 三路交叉 &
MATPOWER 小算例；IEEE 14 节点 &
AC OPF、DC OPF、Newton PF 三者一致性 &
PF 回放收敛；DC 目标 $\le$ AC 目标 $\times 1.05+1$ &
\srcpath{tests/test\_acopf\_dcopf\_crossval.cpp}:250--303；
\srcpath{tests/test\_crossmodule\_integration.cpp}:206--246\\
Native/\allowbreak Parity/\allowbreak Ipopt 后端交叉 &
case9、case30、case300\_acdc、case2000\_acdc &
三后端同案收敛性；Ipopt 有界参数矩阵 &
Parity 目标 $\le$ ED 目标 $\times 1.001+1$；Ipopt 预算内结果有限 &
\srcpath{tests/test\_opf\_solver\_backends.cpp}:1075--1150；
\srcpath{docs/large\_hybrid\_ipm\_diagnostics.md}:118--131、195--220（文档实测）\\
夜间工况 OPF &
南沙全网 JSON、画布 \fld{power\_system.json}（无常规机组） &
外部电网按成本提升为 OPF 变量后收敛 &
\fld{allow\_fallback=false} 下收敛；能量平衡 $10^{-3}$~MW &
\srcpath{tests/test\_nighttime\_opf.cpp}:46--96、164--209\\
时序 UC$\to$OPF$\to$PF 流水线 &
画布系统 24 步/4 步（含 UC） &
逐步 OPF 与 PF 均收敛；UC 可行 &
24/24、4/4 步收敛 &
\srcpath{tests/test\_nighttime\_opf.cpp}:125--162；
\srcpath{docs/time\_series\_power\_flow\_models.md}:276--349\\
三相混合 OPF &
手组 3 母线 9 相节点 + 2 直流节点 &
Full 与 GraphReduced 一致；独立耦合 PF 复现；SOCP 松弛下界 &
目标相对差 $10^{-6}$；电压 $10^{-6}$；VUF $\le 0.05$ &
\srcpath{tests/test\_three\_phase\_hybrid\_opf.cpp}:24--143、284--317、463--510\\
RPO 专项 &
case9；case2869 + 20 OLTC + 10 并联（authored）；case300\_acdc &
独立 OPF 复算 + PF 回放 + 物理网损账 &
目标相对差 $\le\max(10^{-3},10\epsilon_d)$；$|\Delta V|\le\max(5\times10^{-4},10\epsilon_p)$ &
\srcpath{tests/test\_reactive\_power\_opt.cpp}:16--110；
\srcpath{docs/reactive\_power\_optimization\_validation.md}:103--123、159--177（文档实测）\\
大规模混合 IPM 诊断 &
case300\_acdc、case2000\_acdc（2000 节点 + 8 VSC）、2869/9241/13659 &
UMFPACK/\allowbreak KLU/\allowbreak MUMPS/\allowbreak Eigen 多后端同一收敛；文档实测残差 &
测试：违反 $<10^{-5}$、迭代上限；文档：$3.49\times10^{-9}$/$9.03\times10^{-7}$ &
\srcpath{tests/test\_opf\_solver\_backends.cpp}:86--112、880--943、1132--1196；
\srcpath{docs/large\_hybrid\_ipm\_diagnostics.md}:77--131、224--289（文档实测）\\
BPA/DSP LCC 算例 &
\srcpath{2DC.dat}、\srcpath{cigre.dat}（6 交流母线、2 直流母线、2 台 LCC） &
DSP SOL 外部参考；解析 Jacobian/Hessian 有限差分；固定 tap PF 回放 &
$1500/1410$~MW $\pm2$；线损 90~MW $\pm2$；Jacobian $2\times10^{-5}$；回放 $0.5$~MW &
\srcpath{tests/test\_opf\_solver\_backends.cpp}:310--422、424--598、600--878；
\srcpath{docs/lcc\_dat\_opf\_report/LCC\_DAT\_OPF\_technical\_report.tex}:1039--1124\\
结果审计、回放与投影 &
case9、case14、IEEE14 AC/DC、4 母线串联归约 &
\srcpath{verify\_opf\_result} 审计；归约/投影前后目标一致 &
电压违反 $<0.02$~pu；归约目标差 $10^{-3}$ &
\srcpath{tests/test\_solver\_capabilities.cpp}:194--239；
\srcpath{tests/test\_graph\_roundtrip.cpp}:223--258、671--705；
\srcpath{tests/test\_multiscale\_comprehensive.cpp}:170--229\\
\end{longtable}
\end{footnotesize}

\subsection{手算解析解基准}\label{subsec:ver:hand}

\srcpath{external\_data/opf\_hand\_cases/} 提供四个刻意做到极小的 JSON 算例，每个都在
README 中给出完整闭式推导，用于对照 GUI/API 的 DC OPF 路径（直流模型无损，手算结果
精确；README 同时声明 AC/parity 路径只作冒烟，因为交流可行性细节不在手算范围内，
\srcpath{external\_data/opf\_hand\_cases/README.md}:1--19）。

\subsubsection*{算例与校验量}
\begin{itemize}
  \item \textbf{01 单母线 merit order}：60~MW 负荷，机组 1（0--40~MW，成本 10）与
        机组 2（0--100~MW，成本 20）。手算调度 $P_1=40$、$P_2=20$~MW，目标 800，
        单节点 LMP 为边际机组成本 20
        （\srcpath{external\_data/opf\_hand\_cases/README.md}:21--64）；
  \item \textbf{02 两母线支路拥塞}：100~MW 负荷，线路 $x=0.1$~pu、限值 40~MVA。
        手算 $P_1=40$、$P_2=60$~MW，目标 3400；DC 角度按
        \begin{align*}
        f_{12}=\frac{\theta_1-\theta_2}{x_{12}}\,S_{\mathrm{base}}
        \end{align*}
        得 $\theta_2=-0.04$~rad$=-2.2918^\circ$，两侧 LMP 分别为 10 与 50
        （同文件:65--124）；
  \item \textbf{03 负荷削减}：80~MW 负荷、50~MW 机组上限。默认 \fld{voll=0} 时求解器
        取有效 VOLL $=\max(100,\,1.1\,c_1)=100$，手算削减 30~MW、目标 4000
        （同文件:126--174）；
  \item \textbf{04 二次成本的 GUI DC 口径}：README 明确当前 GUI DC 端点对二次成本按
        保守边际系数的 LP 线性化处理，
        \begin{align*}
        c_i^{LP}=c_{1i}+2c_{2i}P_i^{\max},
        \end{align*}
        手算 LP 调度 $(80,0)$~MW、目标 1120；真 QP 参考解 $(36.667,43.333)$~MW、
        目标 815.333 仅供支持真 QP 的后端对照（同文件:176--296）。
\end{itemize}

\subsubsection*{已知偏差}
该目录在当前检出中\textbf{未被任何 C++ 测试或 ctest 目标引用}（全仓检索仅
\srcpath{ngnlp/chapters/07\_validation\_and\_usage.tex}:57--60 提及），定位为
GUI/API 人工校核材料；算例 02 的 README 明确“支路对偶当前不被 DC OPF 路径认证，
应校验原始调度与潮流限值而非 \fld{branch\_mu\_*}”
（\srcpath{external\_data/opf\_hand\_cases/README.md}:124），与
\ref{subsec:ver:limitations} 节的支路对偶条目一致。

\subsection{MATPOWER 算例通道}\label{subsec:ver:matpower}

\subsubsection*{小算例（$\le 57$ 节点）}
\srcpath{tests/test\_acopf\_dcopf\_crossval.cpp}:52--54 固定 case5/9/14/30/39/57
六个算例，AC OPF 走 NativeAC（\fld{use\_parity\_ipm=false}，同文件:117--130），
DC OPF 走自研 NativeQP LCQP（同文件:132--144），两者均以可行性容差 $10^{-6}$ 调用；
随后把 AC OPF 的 $P_g/Q_g/V_g$ 写回系统副本跑 Newton 潮流（容差 $10^{-8}$，
同文件:151--183）。\srcpath{tests/test\_matpower\_cases.cpp}:618--638 另有两项
DC OPF 断言：case9 目标有限且 \fld{pg\_mw} 尺寸正确，case14 所有在役机组调度落在
$[P_{\min}-10^{-3},\,P_{\max}+10^{-3}]$~MW 内。

\subsubsection*{大算例（Parity IPM）}
\srcpath{tests/test\_opf\_solver\_backends.cpp}:1058--1071 是 case2383wp（波兰冬季
峰荷）的回归守卫：注释记录该算例在引入基于梯度的目标缩放前发散（迭代顶格时最大
平稳性约 0.45），现断言收敛。同文件:1154--1176 断言 case2869pegase 收敛、迭代
$<100$ 且 KKT 维数超 1500 时实际使用稀疏分解后端（\fld{linear\_solver\_backend}
含 \fld{sparse}）。同文件:1178--1196 对 case6515rte 只断言目标有限且
$<10^{30}$——注释明确“该难例在迭代预算内不期望收敛”，这是对此前 Eigen SparseLU
下发散到 NaN（目标约 $-3\times10^{48}$）的稳健性守卫，不是收敛性声明。
文档另记录 case9241/13659 的参考目标 315837/386107 与 KLU 后端实测
（\srcpath{docs/large\_hybrid\_ipm\_diagnostics.md}:263--269，文档实测）。
诊断目标 \srcpath{opf\_numerical\_benchmark} 另对 25,000-bus ACTIVSg25k 执行
一次预热加三次计时的 Release/KLU 协议：普通 AC-PF 启动中位数为
$183.808\,\mathrm{s}$，66 次迭代、65 次分解，primal/dual 为
$3.239\times10^{-7}/1.004\times10^{-7}$；默认 $2\,\mathrm{s}$ DC Phase I
未形成可接纳候选，精确回退后中位数为 $186.588\,\mathrm{s}$，慢 $1.51\%$。
这是可求解性和预算失败语义的性能验证，不是 25k 加速声明；该长时诊断不注册到
常规 CTest。
同一目标的实验性零分解 component dispatch-dual predictor 仅把 normalized dual
$468.10733\to468.08762$（$0.0042\%$），未通过固定 $5\%$ exact residual 审计，
以 $1.978\,\mathrm{ms}$ 成本回退；单次总时长 $184.344\,\mathrm{s}$ 且 Phase II
仍为 66 次迭代、65 次分解。因此它默认关闭，结果支持优先研究完整 continuation
state 与跨重复求解 symbolic pattern 复用，而非 partial dual 或更精确的辅助 DC-IPM。
同一 ACTIVSg25k 的 prepared-session 两轮 Release/KLU 验证进一步给出：首轮
$184.645\,\mathrm{s}$/66 次迭代/65 次分解，第二轮
$8.666\,\mathrm{s}$/4/3，且 0 次 symbolic analyze；最终 primal/dual 为
$3.389\times10^{-8}/2.092\times10^{-8}$，峰值 RSS $4709.6\,\mathrm{MiB}$。
这是完整状态与 symbolic reuse 的重复求解收益；第二轮 3 次分解均为新 numeric
factorization，不是旧数值因子的无审计复用。

\subsubsection*{已知偏差}
综合对比用例（\srcpath{tests/test\_acopf\_dcopf\_crossval.cpp}:554--694）跳过全部
$>1000$ 节点的算例（同文件:629--633）与 scenario/contab 文件（同文件:569--574），
且只统计收敛计数（\fld{acopf\_pass>0}），不对目标值与参考解逐一相等断言；
case57 的 PF(DCOPF) 回放收敛被降级为 \texttt{WARN} 而非断言（同文件:273--282，
理由：DC OPF 不提供电压与无功，难例上不保证交流可行）。

\subsection{跨求解器交叉验证}\label{subsec:ver:crossval}

\subsubsection*{AC OPF vs DC OPF vs Newton PF}
\srcpath{tests/test\_acopf\_dcopf\_crossval.cpp}:250--303 对六个小算例逐一断言：
AC OPF 收敛则 PF(AC OPF 调度） 必收敛（同文件:268--271）；DC OPF 目标必须为正
（同文件:287--291）。跨模块流水线测试进一步断言松弛性质
\begin{align*}
C_{\mathrm{DCOPF}} \le 1.05\,C_{\mathrm{ACOPF}} + 1
\end{align*}
（\srcpath{tests/test\_crossmodule\_integration.cpp}:206--209，判据 P1-A4）、
OPF 调度下 Newton PF 必收敛（P1-A6，同文件:232）以及逐母线电压一致
$|V_m^{PF}-V_m^{OPF}|<0.05$~pu（P1-A7，同文件:234--246）。

\subsubsection*{Native / Parity / Ipopt 三路}
\srcpath{tests/test\_opf\_solver\_backends.cpp}:1075--1112 在 case30 上分别断言
EconomicDispatch（merit-order + PF 捷径）、ParityIPM（全空间非线性 OPF）两后端收敛；
同文件:1114--1130 断言真 OPF 目标不超过经济调度捷径的 $1.001$ 倍加 1。Ipopt 后端的
断言是\textbf{优雅降级}：注释声明嵌入式 Ipopt/MUMPS 默认被门控关闭，选择 Ipopt 时必须
回退到 Parity IPM 且 \fld{solver\_path} 报告为 \fld{ParityIPM}（同文件:1100--1108；
编译期门控宏 \srcpath{HACDCPF\_HAVE\_IPOPT}，\srcpath{src/optimal\_power\_flow/ac\_opf.cpp}:33；
CMake 选项 \srcpath{HACDCPF\_ENABLE\_IPOPT} 默认 OFF、macOS 默认 ON，见
第~\ref{sec:overview}~节）。case300\_acdc 上另断言 Ipopt 迭代/容差选项透传生效：
1 次迭代预算必须报 \texttt{Max iterations exceeded} 且不收敛，800 次预算收敛且
违反 $<10^{-6}$（同文件:114--149）。

\subsubsection*{已知偏差}
\begin{itemize}
  \item AC 与 DC 目标值\textbf{不做直接相等断言}：测试注释说明两者成本口径不同，
        直接比较无意义（\srcpath{tests/test\_acopf\_dcopf\_crossval.cpp}:284--291）；
  \item AC OPF 结果不暴露支路潮流，$P_f$ 对比代码留空（同文件:222--227）；
  \item “DC OPF 调度质量”判据宽松：相对调度差只要求 $<1.0$（即 100\%，
        同文件:536--547；注意 546 行尾注释“Less than 50\%”与实际断言 \fld{rel\_diff<1.0}
        不一致，以断言为准）；
  \item 默认构建下 Ipopt 选择静默回退 Parity IPM，因此默认 ctest 中并不存在真正的
        三路目标值相等断言；Ipopt 侧的数值证据来自隐藏性能用例与文档实测
        （\srcpath{docs/large\_hybrid\_ipm\_diagnostics.md}:195--220）；
  \item 独立 \srcpath{power\_models} 模块的 OPF 交叉验证目前主要体现在
        三相混合 OPF 的 NativeIPM/Ipopt 适配器路径（见 \ref{subsec:ver:threephase} 节），
        AC OPF 的 AML builder 无独立对照测试；其源码等价模型见
        \srcpath{docs/modules/power\_models/power\_models\_manual.tex}。
\end{itemize}

\subsection{夜间工况 OPF 与外部电网提升}\label{subsec:ver:nighttime}

\subsubsection*{回归起源}
\srcpath{tests/test\_nighttime\_opf.cpp}:1--5 的文件头注释记录了本通道的回归起源：
修复前 \srcpath{solve\_ac\_opf} 在 \fld{sys.ac.generators} 为空（$n_g=0$）时立即返回
失败；修复后带正成本的外部电网被提升为 OPF 发电机变量，夜间（无光伏、无常规机组）
工况得以收敛。

\subsubsection*{算例与判据}
\begin{itemize}
  \item \textbf{画布 \fld{power\_system.json}}（无常规机组、1 外部电网、1 双绕组
        变压器、含 VSC/DC-DC/储能/直流光伏）：Parity IPM、禁 fallback 下断言收敛；
        储能功率与 authored 设定一致（$2\times10^{-4}$~MW），VSC 无功跟随设定
        （$5\times10^{-2}$~MVAr），外部电网无功 $0<Q<0.20$~MVAr，DC-DC 效率与
        等值电阻的直流平衡按公式逐项核对（$2\times10^{-4}$~MW），交流供给与负荷
        平衡到 $10^{-3}$~MW（\srcpath{tests/test\_nighttime\_opf.cpp}:46--96）；
  \item \textbf{南沙全网}（\fld{nansha\_full\_network.json}，无常规机组、外部电网带
        成本）：夜间用例先断言前提（机组为空、至少一个外部电网成本为正），再断言
        \fld{allow\_fallback=false} 下收敛；昼间对照用例同样收敛（同文件:164--209）；
  \item \textbf{同母线独立调度}：50~MW 负荷母线上同时挂 100 成本机组与 10 成本外部
        电网，断言机组 $P_g\approx0$、外部电网承担全部 50~MW（$10^{-2}$~MW），
        AC 与 DC OPF 各一项（同文件:211--255、312--350）；
  \item \textbf{投影正确性}：死岛剥离后发电机结果保持 authored 顺序（同文件:257--310）；
        闭合开关合并母线后外部电网/机组/负荷全部保留并正确调度（同文件:352--409）。
\end{itemize}

\subsubsection*{已知偏差}
本通道覆盖 2 个工程 JSON 算例与 3 个手工微系统，均为“收敛 + 少量物理量”断言，
没有外部工具对照；外部电网的无功调节能力、成本曲线形态在大系统上的行为未单独验证。

\subsection{时序 UC$\to$OPF$\to$PF 流水线校核}\label{subsec:ver:tspf}

时序生产模拟的三阶段流水线（UC MILP $\to$ 逐步 AC OPF $\to$ 逐步混合潮流回放）
模型细节属第~\ref{sec:overview}~节范围，此处只列与 OPF 可信性直接相关的校核机制
（\srcpath{docs/time\_series\_power\_flow\_models.md}:276--349）：

\begin{itemize}
  \item \textbf{跟踪带}：对每台在线机组，OPF 功率盒围绕 UC 基点收紧，
        \begin{align*}
        \beta_g=\max\bigl(\beta^{abs},\ \beta^{rel}\max(|p^{UC}_g|,1)\bigr),
        \end{align*}
        松弛机 $(\beta^{rel},\beta^{abs})=(20\%,5\,\mathrm{MW})$，非松弛机
        $(2\%,1\,\mathrm{MW})$；其工程语义是 OPF 只作类 AGC 再调度，吸收无损 UC
        看不见的网损与直流近似误差（\srcpath{docs/time\_series\_power\_flow\_models.md}:284--300）；
  \item \textbf{降级路径}：某步 OPF 不收敛时，该步退回对原始 UC 快照直接跑 PF，
        流水线不中断（同文件:307--308）；
  \item \textbf{交叉指标}：每步记录 $\max_b|V_{m,b}^{OPF}-V_{m,b}^{PF}|$ 与经首母线
        参考对齐的相角差、双侧物理网损，作为流水线“信任证书”（同文件:323--339）；
  \item \textbf{储能一致性}：OPF 故意不重调度储能，避免功率与 SOC 轨迹互相矛盾
        （同文件:302--305）；
  \item \textbf{线程安全}：按日并行时逐步 OPF 被强制到线程安全的 Parity IPM，
        不可重入的 Ipopt/MUMPS 路径不会并发进入（同文件:412--414）。
\end{itemize}

测试侧：\srcpath{tests/test\_nighttime\_opf.cpp}:125--162 断言画布系统 24 步动态
OPF 全部 24 步 OPF 与 PF 收敛，以及含 UC 的 4 步默认日生产路径 UC 可行且
4/4 步收敛。

\subsection{三相混合 OPF 测试}\label{subsec:ver:threephase}

模型与算法见第~\ref{sec:threephase}~节；本小节只列验证证据。全部用例共用同一个
手组算例：3 个三相母线（9 个相节点）、2 个直流节点、1 台 VSC（效率 0.98、容量
0.25~pu）、3 台相机组，电压不平衡度上限 \fld{vuf\_max}$=0.05$
（\srcpath{tests/test\_three\_phase\_hybrid\_opf.cpp}:24--102）。

\subsubsection*{验证项与容差}
\begin{itemize}
  \item \textbf{Full vs GraphReduced}：两种模型变体在同一算例上均收敛，目标相对差
        $\le10^{-6}$，全电压复数差 $<10^{-6}$，电压违反 $<10^{-8}$，VUF 不超过
        上限 $+10^{-7}$，换流器电流不平衡度一致到 $10^{-6}$
        （同文件:106--143）；
  \item \textbf{图归约序列}：两断面序列求解与独立求解一致（目标 $10^{-7}$、电压
        $10^{-6}$），归约不丢时变负荷节点（同文件:145--172）；
  \item \textbf{GFL/GFM 动态平衡}：在 OPF 运行点上初始化 GFL（PLL）与 GFM（下垂）
        动态模型并 trim 到网络平衡，状态导数 $<10^{-7}$；解析等式 Jacobian 与
        Lagrangian Hessian 误差均 $<10^{-5}$（同文件:174--282）；
  \item \textbf{独立耦合 PF 复现}：由 OPF 结果构造三相混合 PF 算例，独立 Newton
        PF 复现电压到 $2\times10^{-6}$（同文件:284--317）；
  \item \textbf{Native Full 认证归约解}：把 GraphReduced 的原始--对偶点传输到 Full
        模型热启动，初始/最终原始与对偶残差、互补均 $\le10^{-6}$，目标一致到
        $10^{-6}$（同文件:319--401）；
  \item \textbf{约束 oracle}：惰性不等式求解与全行求解目标一致到 $10^{-7}$，
        被略去行的裕度满足激活阈值（同文件:403--461）；
  \item \textbf{SOCP 松弛下界}：抬升松弛给出经对偶证书认证的目标下界，不超过
        非线性解 $+10^{-7}$，外逼近轮次下界单调不降（同文件:463--510）。
\end{itemize}

\subsubsection*{已知偏差}
算例规模极小（9 相节点），不构成大规模三相系统的性能证据；松弛路径对恒电流负荷
显式拒绝并报告 \fld{model\_limitations}（同文件:512--522）；测试使用 AML 引擎层的
\fld{SolverBackend::Ipopt}/\fld{NativeIPM}，其可用性随 MIPSolvers 后端配置变化。


\subsection{RPO 专项验证}\label{subsec:ver:rpo}

模型（离散调压设备搜索 + 连续 AC/DC OPF 补全）见第~\ref{sec:rpo}~节；验证证据分
测试断言与文档实测两部分。

\subsubsection*{测试断言（\srcpath{tests/test\_reactive\_power\_opt.cpp}）}
\begin{itemize}
  \item \textbf{诚实局部证书}：case9 网损目标 RPO 收敛，且结果明确
        \fld{globally\_certified=false}、\fld{optimality\_gap\_available=false}、
        \fld{model\_limitations} 非空，算法标识为
        \fld{discrete\_coordinate\_search\_with\_ac\_opf}；无可调离散设备时目标
        等于基线网损（同文件:16--42）；
  \item \textbf{可投切并联}：case9 加 2 步电容器组后，电压偏差目标不差于基线
        （$+10^{-10}$），档位动作在铭牌范围内；\fld{gap} 恒为 1.0（不报告虚假
        全局间隙，同文件:44--82）；
  \item \textbf{控制清单}：\srcpath{inspect\_rpo\_controls} 对可调/固定/档位越界/
        停运设备分别给出排除原因，档位窗口 \fld{max\_tap\_move} 生效（同文件:112--182）；
  \item \textbf{MATPOWER 固定变比不伪造}：MATPOWER 导入的支路变比保持固定，
        不再虚构 4001 档 OLTC；\srcpath{case300\_acdc} 仅 authored 3 台
        $[-4,4]$、每档 1.25\% 的 OLTC 可调（同文件:184--211）；
  \item \textbf{档位动作联动}：带 \fld{source\_branch\_idx} 的 OLTC 动作按相对变比
        同步更新关联 AC 支路 \fld{tap}，保证 RPO 内层、独立 OPF 与 PF 回放使用
        同一调压动作（同文件:213--243）；
  \item case300 种子求解用例带 \fld{[.performance]} 标签，默认不运行（同文件:84--110）。
\end{itemize}

\subsubsection*{全过程交叉验证（文档口径）}
每次 RPO 后固定执行“保存离散档位 $\to$ 独立 OPF 复算 $\to$ PF 回放 $\to$ 网损账”
四步（\srcpath{docs/reactive\_power\_optimization\_validation.md}:103--123）。
独立 OPF 一致性判据为
\begin{align*}
\frac{|f_1-f_2|}{\max(1,|f_1|)}\le\max(10^{-3},10\epsilon_d),\qquad
\max_b|\Delta V_b|\le\max(5\times10^{-4},10\epsilon_p),
\end{align*}
其中 $\epsilon_d$、$\epsilon_p$ 分别为平稳性与可行性容差；非凸 OPF 可能存在目标
等价但 P/Q 分配不同的局部点，故 P/Q 最大差异只展示、不触发假失败。纯 AC 的 PF
电压差判据为 $10^{-3}$~pu；富混合 AC/DC 在 authored-space 调度未完整回写时显示
“可用数值检查通过，PF 回放部分覆盖”（同文件:113--123）。

\subsubsection*{大规模实测（文档实测）}
case2869 + 20 台 OLTC + 10 台并联的 authored 系统上：单位成本经济基线 74 次迭代
/25~s 收敛到参考目标 133999（违反 $1.1\times10^{-6}$）；暖启动链 2--7 次迭代、
平均约 0.46~s/次；完整 MinActiveLoss RPO 动作 22 台设备，物理网损
$1561.8\to1558.8$~MW。IEEE24 三区域回归中 vdev（MatchRPO 路径）与旧版逐位一致；
loss 路径因改用等价经济基点，基线/最终网损均优于旧版，combined 正常收敛
（\srcpath{docs/reactive\_power\_optimization\_validation.md}:159--177）。

\subsubsection*{已知偏差}
稳健内层 \fld{LossEconomic} 按网损最优分配机组无功、不直接为电压最优：case2869 上
vdev RPO 虽收敛，但 vdev $0.1392\to0.1401$ 基本未降；3000+ 完整搜索深度受时间限
（默认 120~s 内得可行改进点），离散评估并行化未完成
（\srcpath{docs/reactive\_power\_optimization\_validation.md}:179--191）。

\subsection{大规模混合 IPM 诊断证据}\label{subsec:ver:ipm}

\subsubsection*{case300\_acdc 根因与修复}
\srcpath{docs/large\_hybrid\_ipm\_diagnostics.md}:45--75 记录：故障并非模型不可行
（Ipopt 可把原始残差压到 $2.5\times10^{-9}$），而是原生 IPM 全步更新缺全局化，
轨迹进入病态区后所有稀疏分解失败。修复是对维数 $\ge512$ 的 KKT 增加残差 filter
回溯（原始可行性/对偶平稳性/互补三轴，步长减半最多 14 次），位于
\srcpath{src/optimal\_power\_flow/parity\_ipm.cpp}。修复后 case300\_acdc 原生 IPM
42 次迭代收敛，最大约束残差约 $1.74\times10^{-7}$、平稳性约 $2.77\times10^{-7}$，
41 接受步/17 回溯拒绝；强制 UMFPACK、KLU、Eigen SparseLU、Dense LU 四后端均通过
同一算例（同文件:77--88，文档实测）。测试侧对应断言为
\srcpath{tests/test\_opf\_solver\_backends.cpp}:86--112（收敛、接受步 $>0$ 且
$\le$ 迭代数、拒绝步 $>0$）。

\subsubsection*{2000 节点富混合实测}
\srcpath{docs/large\_hybrid\_ipm\_diagnostics.md}:118--131 记录（文档实测，Release）：
原生 parity IPM + 稀疏 UMFPACK 90 次迭代/7.07~s 收敛，最大违反
$3.49\times10^{-9}$、平稳性 $9.03\times10^{-7}$；内嵌 Ipopt + MUMPS 100 次预算
48.15~s 后原始残差 $4.96\times10^{-4}$、对偶 1.49，未达停止判据。测试断言为
\srcpath{tests/test\_opf\_solver\_backends.cpp}:880--912（Parity IPM 收敛、
违反 $<10^{-5}$、迭代 $<160$）与同文件:914--943（Ipopt 有界基准，仅断言有限，
隐藏性能用例）。据此文档建议 2000 节点级富混合系统优先使用原生 IPM，Ipopt 保留为
小中型问题后端与交叉诊断路径。

\subsubsection*{后端结构分工与暖启动}
2026-07-19 更新的诊断结论（同文件:224--289，文档实测）：线性代数默认顺序
MUMPS $\to$ UMFPACK $\to$ KLU $\to$ Eigen；Newton 系统默认增广 KKT +
W\"achter--Biegler $\delta_W$ 惯性校正；纯 AC 大算例自动走 KLU（BTF 预分解契合
近块三角电网 KKT，2869/9241/13659 三案收敛到同一参考目标且更快），混合 AC/DC 因
增广 KKT 惯性信息只有 MUMPS 报告而必须走 MUMPS——文档明确 KLU/LU 不报告惯性时
混合算例失去正则而停滞。\fld{ac\_pf\_warm\_start} 选项使 case13659pegase 以 189 次
迭代/约 41~s 收敛于参考目标 386107，但在良态算例上反而更慢、个别 rte 算例不收敛，
因此是选择性启用而非默认（同文件:235--239）。

\subsubsection*{已知偏差}
arm64 本地静态 MUMPS 5.7.3 在全部六类 NLP 测试中第 1 次迭代即
\texttt{Ipopt restoration failed}，当前默认仍用 Homebrew MUMPS 动态库
（同文件:111--116、190--193）；Jacobian 系数尺度达 $10^{12}$ 时默认 gradient
scaling 会压小 scaled 对偶误差，适配器在启用未缩放分量门槛后明确返回
\texttt{search direction too small} 而非假收敛——此类模型必须在建模层无量纲化
（同文件:184--188）。

\subsection{BPA/DSP LCC 算例验证}\label{subsec:ver:lcc}

LCC 直流算例的 OPF 验证是平台唯一带\textbf{外部商业软件参考解}的 OPF 通道。
完整证据链见专项报告
\srcpath{docs/lcc\_dat\_opf\_report/LCC\_DAT\_OPF\_technical\_report.tex}，
此处汇总与本手册相关的部分。

\subsubsection*{算例与基准}
\srcpath{data/dsp/cigre.dat}（DSP Samples/CIGRE）：100~MVA 基准，6 交流母线、
4 交流支路、2 直流母线、1 条 $10~\Omega$ 直流线路、2 台 LCC；DSP 参考解为
整流/逆变 $U_d=500/470$~kV、$I_d=3$~kA、交流有功 $-1500/+1410$~MW、线损
$I_d^2R=90$~MW（同报告 :982--1009）。\srcpath{data/dsp/2DC.dat} 为补充交叉校核算例
（同报告 :1116--1124）。

\subsubsection*{验证项与容差}
\begin{itemize}
  \item \textbf{导数级}：LCC parity 平衡方程解析 Jacobian 对中心有限差分
        $<2\times10^{-5}$，覆盖阀侧母线、换相母线与直流电压三类变量
        （\srcpath{tests/test\_opf\_solver\_backends.cpp}:310--422）；隔离 LCC 的
        局部 Hessian 对称性 $<10^{-10}$，对 Jacobian 差分相对容差 $5\times10^{-3}$
        （同文件:424--598）；
  \item \textbf{raw CIGRE OPF}：原始 DAT 直接导入（BS 平衡机 $P_{\min}=-9999$~MW
        为卡片语义），Parity IPM 收敛，$P_g$ 复现 $1500/-1410$~MW $\pm2$、线损
        90~MW $\pm2$、违反 $<10^{-6}$；OPF 结果 JSON 双向往返不丢
        \fld{lcc\_transfers}（同文件:600--671）；报告记录 7 次迭代收敛、
        primal/dual 残差 $1.7\times10^{-11}$/$6.0\times10^{-7}$（同报告 :1095--1108，
        文档实测）；
  \item \textbf{不可行变体安全失败}：受端 $P_{\min}=0$ 时模型真不可行，断言返回
        有限、可序列化的失败结果，不伪造无损直流解、不从错误尺寸 primal 提取物理
        运行点（同文件:673--710）；
  \item \textbf{固定 tap PF 回放}：把 OPF 调度写回副本（显式关闭 PF 专用 R 卡
        分接头外环以保证同一 Ybus），LCC 端口 $P/Q/P_{dc}$ 与 OPF 一致到
        0.5~MW/MVAr（同文件:814--846）；
  \item \textbf{后端守卫}：EconomicDispatch 后端对混合 AC/DC/LCC 显式拒绝；
        DC OPF 显式拒绝并打 \fld{dc-opf:rejected-lcc} 范围标签；不支持的控制模式
        组合与缺失换相电抗在 readiness 检查即拒绝（同文件:848--877）。
\end{itemize}

\subsubsection*{当前验证状态（报告声明）}
报告“当前验证状态”表声明：DAT 接口与 PF、DSP 对比、LCC OPF（175 项断言/4 个
用例，含三变量 Jacobian/Hessian、raw CIGRE、固定 tap 回放与后端保护）、
$P_{\min}=0$ 不可行变体、浏览器 CIGRE E2E 均通过
（\srcpath{docs/lcc\_dat\_opf\_report/LCC\_DAT\_OPF\_technical\_report.tex}:1039--1060）。
同一节明确\textbf{完整 \srcpath{test\_opf\_solver\_backends} 不在该轮验收范围内}，
且历史全量运行存在与 LCC 无关的后端/大算例基线失败，不能写成全绿（同文件
:1058--1060）。

\subsubsection*{已知偏差}
OPF 中换流变压器 tap 是固定 formulation 输入，不复现 DSP 的调档运行点；报告接受
准则也不要求复现调档点（同文件:1036--1037、1086--1114）。2DC 卡片
$P_{\min}=P_{\max}=0$，是 PF/Jacobian 参考而非 OPF-ready（测试同文件:712--715）。
其余局限（准稳态、无换相重叠迭代、clamp 不可二阶光滑等）见
\ref{subsec:ver:limitations} 节。

\subsection{结果审计、PF 回放与投影 round-trip}\label{subsec:ver:audit}

\begin{itemize}
  \item \textbf{\srcpath{verify\_opf\_result} 审计}：case9 收敛结果经审计后
        电压限值违反 $<0.02$~pu、机组限值违反 $<1.0$~MW，报告/重算目标均有限；
        结构性不可行系统（负荷无源）审计必须给出不可行提示
        （\srcpath{tests/test\_solver\_capabilities.cpp}:194--239）；
  \item \textbf{混合系统禁 AC-only 回退}：含 DC 母线与 VSC 的系统在极端容差下
        不收敛时，状态不得含 \texttt{economic-dispatch + AC PF fallback}，且不可行
        提示中必须出现 \texttt{AC-only economic-dispatch fallback suppressed}
        （同文件:241--302）；
  \item \textbf{能力标志与实际行为一致}：\fld{supports\_ac\_opf}/
        \fld{supports\_dc\_opf} 恒真且实解收敛（同文件:135--156）；
  \item \textbf{投影 round-trip}：纯 AC 重编号后 PF 与 AC OPF 结果向量尺寸与
        authored 母线一致（\srcpath{tests/test\_graph\_roundtrip.cpp}:223--258）；
        串联归约前后 DC OPF 目标一致到 $10^{-3}$（同文件:671--705）；
  \item \textbf{AML 数据装配}：孤立 DC 母线及其关联支路不进入 OPF 模型
        （\srcpath{tests/test\_hacdcpf.cpp}:234--261）；
  \item \textbf{全约束族冒烟}：多尺度综合算例在支路限值、换流器容量/电流/调制
        约束全部启用下收敛，违反 $<100\times$ 可行性容差
        （\srcpath{tests/test\_multiscale\_comprehensive.cpp}:170--229）；
  \item \textbf{入口冒烟}：IEEE14 AC/DC 的 AC/DC OPF 入口返回非空状态串
        （\srcpath{tests/test\_hacdcpf.cpp}:372--400）；
  \item \textbf{浏览器端到端}：RPO GUI 全流程（输入、运行、结果与诚实证书）由
        Playwright E2E 覆盖，注册于 \srcpath{tests/CMakeLists.txt}:532--537。
\end{itemize}


\subsection{已知局限汇总}\label{subsec:ver:limitations}

表~\ref{tab:ver:limitations} 从测试中的 skip/WARN 声明、门控分支、结果结构中的
诚实标签与文档“局限/已知问题”段落逐条收集，未做美化。表中“表现/依据”列为代码
或文档原文语义的转述。

\begin{footnotesize}
\begin{longtable}{@{}L{4.2cm}L{1.9cm}L{5.2cm}L{4.0cm}@{}}
\caption{OPF 模块已知局限汇总}\label{tab:ver:limitations}\\
\toprule
\textbf{局限} & \textbf{影响范围} & \textbf{表现/依据} & \textbf{出处（文件:行号）}\\
\midrule
\endfirsthead
\multicolumn{4}{@{}l}{\footnotesize（续表）}\\
\toprule
\textbf{局限} & \textbf{影响范围} & \textbf{表现/依据} & \textbf{出处（文件:行号）}\\
\midrule
\endhead
\midrule
\multicolumn{4}{r@{}}{\footnotesize（接下页）}\\
\endfoot
\bottomrule
\endlastfoot
嵌入式 Ipopt 默认被门控关闭 &
AC OPF 后端 &
\srcpath{HACDCPF\_ENABLE\_IPOPT} 默认 OFF（macOS ON）；选 Ipopt 而后端不可用时静默回退 Parity IPM，\fld{solver\_path} 报告 \fld{ParityIPM} &
\srcpath{tests/test\_opf\_solver\_backends.cpp}:1100--1108；\srcpath{src/optimal\_power\_flow/ac\_opf.cpp}:33\\
Ipopt 在 2000 节点混合算例预算内不达标 &
大规模混合 &
100 次预算后原始 $4.96\times10^{-4}$、对偶 1.49，未收敛；文档建议生产路径用原生 IPM &
\srcpath{docs/large\_hybrid\_ipm\_diagnostics.md}:118--131、211--220（文档实测）\\
case6515rte 不期望收敛 &
大 AC 算例 &
测试仅断言目标有限（$<10^{30}$）与稀疏后端，注释明确不期望预算内收敛 &
\srcpath{tests/test\_opf\_solver\_backends.cpp}:1178--1196\\
Ipopt 有界基准与参数矩阵为隐藏用例 &
回归可见性 &
带 \fld{[.performance]} 标签，默认 ctest 不运行；case300 RPO 种子用例同 &
\srcpath{tests/test\_opf\_solver\_backends.cpp}:914--946；\srcpath{tests/test\_reactive\_power\_opt.cpp}:84--85\\
AC 与 DC OPF 目标值不直接可比 &
交叉验证 &
测试注释声明成本口径不同，直接比较无意义；只断言各自收敛与 DC 目标为正 &
\srcpath{tests/test\_acopf\_dcopf\_crossval.cpp}:284--291\\
AC OPF 结果不暴露支路潮流 &
交叉验证 &
$P_f$ 对比代码留空注释“skip for now” &
\srcpath{tests/test\_acopf\_dcopf\_crossval.cpp}:222--227\\
DC OPF 调度质量判据宽松 &
交叉验证 &
相对调度差断言 $<1.0$（100\%）；546 行尾注释“50\%”与断言不符，以断言为准 &
\srcpath{tests/test\_acopf\_dcopf\_crossval.cpp}:536--547\\
case57 的 PF(DCOPF) 回放不保证收敛 &
交叉验证 &
收敛性降级为 \texttt{WARN}，非断言 &
\srcpath{tests/test\_acopf\_dcopf\_crossval.cpp}:273--282\\
综合对比限于 $\le1000$ 节点 &
交叉验证 &
更大算例 SKIP；scenario/contab 文件排除 &
\srcpath{tests/test\_acopf\_dcopf\_crossval.cpp}:569--574、629--633\\
DC OPF 支路对偶未认证 &
DC OPF LMP &
含 \fld{pf} 变量界时 native simplex 不支持对偶提取，\fld{branch\_mu\_valid=false} &
\srcpath{tests/test\_acopf\_dcopf\_crossval.cpp}:794--799、868--872\\
LMP 提取依赖辅助 LP &
DC OPF LMP &
仅 \fld{include\_branch\_limits=false}（无 \fld{pf} 变量）时辅助 LP 成功 &
\srcpath{tests/test\_acopf\_dcopf\_crossval.cpp}:698--702、716--719\\
Ipopt 乘子不全时 LMP 标不可用 &
AC OPF &
适配器不返回完整约束乘子时 LMP 必须标为不可用 &
\srcpath{docs/lcc\_dat\_opf\_report/LCC\_DAT\_OPF\_technical\_report.tex}:1227--1228\\
混合系统抑制 AC-only 经济调度回退 &
AC OPF &
含 DC/LCC 时禁止丢弃混合网络；提示 \texttt{fallback suppressed} &
\srcpath{tests/test\_solver\_capabilities.cpp}:241--302\\
EconomicDispatch 后端拒绝混合/LCC &
AC OPF &
状态串含 \texttt{hybrid AC/DC/LCC}，不收敛 &
\srcpath{tests/test\_opf\_solver\_backends.cpp}:848--854\\
DC OPF 拒绝 LCC &
DC OPF &
显式失败并打 \fld{dc-opf:rejected-lcc} 范围标签 &
\srcpath{tests/test\_opf\_solver\_backends.cpp}:856--859\\
AML builder 不支持 LCC &
AML 建模层 &
入口/能力范围必须显式声明，不把 LCC 当 0 注入 &
\srcpath{docs/lcc\_dat\_opf\_report/LCC\_DAT\_OPF\_technical\_report.tex}:954\\
三相混合 OPF 不支持在役 LCC &
三相 OPF &
GUI/HTTP 入口显式拒绝并提示改用平衡 AC/DC Parity 或 Ipopt &
\srcpath{docs/lcc\_dat\_opf\_report/LCC\_DAT\_OPF\_technical\_report.tex}:955\\
LCC 换流变 tap 在 OPF 中为固定输入 &
LCC OPF &
结果须声明 \texttt{taps are fixed inputs} 与 \texttt{no tap-changer control}；联合优化需后续把 $t_k$ 纳入决策 &
\srcpath{docs/lcc\_dat\_opf\_report/LCC\_DAT\_OPF\_technical\_report.tex}:1086--1114、1210--1212\\
LCC 为快照准稳态模型 &
LCC &
不含换相重叠角显式迭代、换流失败、谐波、保护与控制动态；不能推断暂态安全 &
\srcpath{docs/lcc\_dat\_opf\_report/LCC\_DAT\_OPF\_technical\_report.tex}:1206--1207\\
LCC 额定电流 clamp 不可二阶光滑 &
LCC OPF &
活动区切换点非二阶光滑；应报告绑定状态，局部 KKT 非全局证书 &
\srcpath{docs/lcc\_dat\_opf\_report/LCC\_DAT\_OPF\_technical\_report.tex}:1213--1214\\
2DC 卡片非 OPF-ready &
LCC 测试 &
$P_{\min}=P_{\max}=0$，仅作 PF/Jacobian 参考 &
\srcpath{tests/test\_opf\_solver\_backends.cpp}:712--715\\
完整 backends 测试目标存在历史基线失败 &
测试套件 &
LCC 报告声明完整 \srcpath{test\_opf\_solver\_backends} 不在该轮验收范围，历史全量有与 LCC 无关的后端/大算例失败 &
\srcpath{docs/lcc\_dat\_opf\_report/LCC\_DAT\_OPF\_technical\_report.tex}:1058--1060\\
RPO 不提供全局最优性证书 &
RPO &
\fld{globally\_certified=false}、\fld{gap=1.0}；局部灵敏度面不用于安全剪枝 &
\srcpath{tests/test\_reactive\_power\_opt.cpp}:16--42、79--81；\srcpath{docs/reactive\_power\_optimization\_validation.md}:82--87\\
RPO 稳健内层不直接为电压最优 &
RPO 大电网 &
case2869 上 vdev RPO 收敛但 vdev 基本未降（$0.1392\to0.1401$） &
\srcpath{docs/reactive\_power\_optimization\_validation.md}:179--191\\
RPO 富混合回放部分覆盖 &
RPO &
混合换流器调度未完整回写 authored-space PF 时标“部分覆盖”；三相单体 RPO 未覆盖 &
\srcpath{docs/reactive\_power\_optimization\_validation.md}:95--99、113--123；\srcpath{docs/large\_hybrid\_ipm\_diagnostics.md}:133--138\\
三相松弛拒绝恒电流负荷 &
三相 OPF &
\fld{fixed-current} 显式拒绝并报告 \fld{model\_limitations} &
\srcpath{tests/test\_three\_phase\_hybrid\_opf.cpp}:512--522\\
手算算例未自动化 &
验证覆盖 &
\srcpath{external\_data/opf\_hand\_cases/} 无 C++ 测试/ctest 引用，仅 GUI 人工校核 &
\srcpath{external\_data/opf\_hand\_cases/README.md}:1--19\\
arm64 静态 MUMPS 数值/ABI 问题 &
Ipopt 后端 &
源码构建 MUMPS 5.7.3 全部用例第 1 次迭代 \texttt{restoration failed}；默认仍用 Homebrew 动态库 &
\srcpath{docs/large\_hybrid\_ipm\_diagnostics.md}:111--116、190--193\\
尺度 $10^{12}$ 模型的缩放假象 &
Ipopt 适配器 &
默认 gradient scaling 压小 scaled 对偶误差；需建模层无量纲化，不能仅靠收紧 \fld{tol} &
\srcpath{docs/large\_hybrid\_ipm\_diagnostics.md}:184--188\\
时序 OPF 单步失败时降级 &
时序流水线 &
该步退回 UC 快照直接 PF，结果精度由交叉指标揭示，不标记为 OPF 质量 &
\srcpath{docs/time\_series\_power\_flow\_models.md}:307--308\\
外部电网夜间通道覆盖有限 &
夜间 OPF &
仅 2 个工程 JSON + 3 个微系统，无外部工具对照 &
\srcpath{tests/test\_nighttime\_opf.cpp}:46--96、164--209\\
\end{longtable}
\end{footnotesize}

\subsection{可信边界结论}\label{subsec:ver:boundary}

综合上述证据，当前 OPF 模块的可信边界可划分为四档。

\subsubsection*{验证充分（默认 ctest 断言 + 外部/文档基准一致）}
\begin{itemize}
  \item \textbf{纯 AC 中小规模经济 OPF}（case5--case300 级）：NativeAC 与 Parity IPM
        均收敛，PF 回放、审计、LMP 符号/量级、拥塞再调度等系统性判据完整；
  \item \textbf{纯 AC 大规模 Parity IPM}：case2383wp、case2869pegase 有默认断言；
        文档实测 case2869/9241/13659 收敛到参考目标值 133999/315837/386107；
  \item \textbf{交直流混合 Parity IPM}：case300\_acdc（四种线性代数后端）与
        case2000\_acdc（2000 节点 + 多端直流 + 8 VSC）有默认断言与文档实测残差；
        暖启动、filter 回溯、稀疏 KKT 路径均有针对断言；有界 Phase I 在
        小型 Hybrid 上认证 primal/dual，在 case300/case2000 严重冷启动上
        只保证单调/预算/诚实非接受，未接受状态不改变 Phase II 起点；
  \item \textbf{DC OPF}：MATPOWER 小算例收敛、限值、孤岛削减、LMP 符号/均匀性/
        拥塞经济性断言完整。
\end{itemize}

\subsubsection*{有测试但覆盖有限（特定算例或单元级断言）}
\begin{itemize}
  \item \textbf{LCC OPF}：2 个 BPA DAT 算例 + DSP 外部参考 + 有限差分导数 + 固定
        tap 回放；多端、定电流/定 $\alpha$ 控制组合、tap 联合优化未覆盖；
  \item \textbf{三相混合 OPF}：9 相节点手组算例上的 Full/归约/松弛/oracle/动态
        平衡全链断言；另有 360 个三相母线、2175 变量的 Phase I 稀疏预算与
        primal/dual-fit 证书，但无同规模完整 Phase II 或外部工具 OPF 对照；
  \item \textbf{RPO}：case9 单元断言 + case2869 authored 文档实测；全局最优性
        明确不声称，大电网电压偏差改善幅度有限；
  \item \textbf{夜间/外部电网 OPF}：2 个工程算例 + 3 个微系统；
  \item \textbf{时序流水线}：24 步动态 OPF 与 4 步 UC 全链断言；多日并行、年度
        8760 步的 OPF 质量由降级路径与交叉指标保障而非逐断言。
\end{itemize}

\subsubsection*{仅有界性/有限性保证}
case6515rte（40 次预算内只断言有限）；case2000\_acdc 的 Ipopt 路径（100 次预算
内文档实测未达停止判据）；case300\_acdc 的 Ipopt 路径需 800 次预算与特定容差组合
方收敛，且平稳性收紧到 $10^{-4}$ 反而退化（文档实测）。这些路径可用于交叉诊断，
不应作为生产精度来源。

\subsubsection*{明确不支持并被守卫}
DC OPF 与 AML builder 的 LCC、三相混合 OPF 的在役 LCC、EconomicDispatch 后端的
混合/LCC 算例、三相松弛的恒电流负荷、混合系统的 AC-only 回退——均在入口或
readiness 检查显式拒绝并打范围标签，不会静默给出误导性结果。

\subsubsection*{总体判断}
OPF 模块在“纯 AC 经济调度”与“平衡 AC/DC 混合 Parity IPM”两条主线上具备从
导数到外部基准的完整证据链，可用于生产；LCC、三相、RPO、时序四条扩展线各有
针对性测试与诚实的能力标签，但其证据是单元/算例级的，使用时应以结果中的
\fld{model\_scope}、\fld{model\_limitations} 与审计字段为准。任何超出上述边界的
规模或工况（$>3000$ 节点富混合、含 LCC 的多端直流优化、三相不平衡单体 RPO），
当前只有工程合理性推断，没有验证证据支撑。
